%%%%%%%%%%%%%%%%%%%%%%%%%%%%%%%%%%%%%%%%%%%%%%%%%%%%%%%%%%%%%%%%%%%%%%
% Câu chưa có mã dạng — cần gắn %[..] rồi chuyển vào đúng file dạng
% Số câu: 3  (sinh tự động bởi tach_cau.py; từ nay thầy sửa trực tiếp file này)
%%%%%%%%%%%%%%%%%%%%%%%%%%%%%%%%%%%%%%%%%%%%%%%%%%%%%%%%%%%%%%%%%%%%%%

%%%%%%%%%%%%%%%% TRẮC NGHIỆM 4 PHƯƠNG ÁN %%%%%%%%%%%%%%%%

% nguồn: Ok_TracNghiemLop11_Lan1
\begin{ex}[c]%[ma:41b6b60b]
	[$\tan a+\tan b=\left(1-\tan  a\cdot\tan b\right)\tan\left(a+b\right)$]%[0D6B3-1]
	Giá trị của biểu thức $\tan 20^\circ+\tan 40^\circ+\sqrt  3\tan 20^\circ\cdot \tan40^\circ $ bằng
	 \choice
		{$-\dfrac{\sqrt  3}{3}$}
		 {$\dfrac{\sqrt 3}{3}$}
		 {$-\sqrt 3 $}
		{\True  $\sqrt 3 $}
	 \loigiai{
		 $
		\begin{aligned}
			\text{Ta  có }&\tan\left(a+b\right)=\dfrac{\tan a+\tan b}{1-\tan a.\tan b}
			\Rightarrow\tan a+\tan b=\left(1-\tan  a\cdot\tan b\right)\tan\left(a+b\right).\\
			\text{Suy ra  }&\tan 20^\circ+\tan 40^\circ=\left(1-\tan 20^\circ\cdot\tan 40^\circ\right)\tan 60^\circ\\
			\Leftrightarrow&\tan 20^\circ+\tan 40^\circ=\sqrt 3-\sqrt  3\tan 20^\circ\cdot\tan 40^\circ\\
			\Leftrightarrow&\tan 20^\circ+\tan 40^\circ+\sqrt 3\tan  20^\circ \cdot\tan 40^\circ=\sqrt 3.
		\end{aligned}
	$} 
\end{ex}

% nguồn: Ok_TracNghiemLop11_Lan1
\begin{ex}[c]%[ma:43f71531]
	[$3^2+4^2=25$]%[0D6K3-1] 
	Cho $M=3\sin  x+4\cos x$. Chọn khẳng định \textbf{đúng}.
	\choice 
		 {$M\le 5$}
		 {$M > 5$}
		 {$M\ge 5$}
		{\True $-5\le M\le 5$} 
	 \loigiai{
		$M=5\left(\dfrac{3}{5}\sin x+\dfrac{4}{5}{\cos x}\right)=5\sin\left(x+a\right)$ với  $\cos a=\dfrac{3}{5},\,\sin a=\dfrac{4}{5}$.\\
	 Ta có $-1\le\sin\left(x+a\right)\le 1\Leftrightarrow-5\le 5\sin\left(x+a\right)\le 5$.}
\end{ex}

% nguồn: Ok_TracNghiemLop11_Lan1
\begin{ex}%[ma:231bea22]
	\immini{ 
	Đường cong trong hình bên là  đồ thị của hàm số nào?%[1D6H3-3]
	 \choice
		{$y=\left(\dfrac{1}{2}\right)^x$} 
		{$y=\log_{\tfrac{2}{5}}  x$}
		{$y=\log_3  x$}
		{\True  $y=2^x$}
		 }{
		\begin{tikzpicture}[>=stealth,scale=0.7,line join=round,line cap=round]
				 \def\xmin{-2} \def\xmax{3} \def\ymax{3}
				\draw[->] (\xmin,0)--(\xmax,0)  node [below]{$x$};
				\draw[->] (0,-0.9)--(0,\ymax) node [left]{$y$}; 
				\node  at (0,0) [below left]{$O$};
				\clip  (\xmin,-0.9) rectangle (\xmax-0.1,\ymax-0.1);
				\draw[smooth,samples=300] plot  (\x,{2^(\x)});
			\end{tikzpicture}
		} 
	\loigiai{ 
		Quan sát hình dáng đồ thị ta thấy đồ thị nằm trên trục $Ox$, đồ  thị đi lên từ trái sang phải nên hàm số đồng biến trên $\mathbb{R}$, vậy chỉ có thể là đồ thị của hàm số $y=2^x$.
	} 
\end{ex}
